\section{Minimal enclosure and the final limits}
\label{ch:final:section}

We now complete the proof of Theorem~\ref{ch:intro:main}. All geometric data and preparation choices are fixed until the last step.

\begin{lemma}\label{ch:final:minimal-enclosure}
Let $(X,\widetilde g)$ be a comparison exterior furnished by the preceding sections. It has a smooth compact nonempty minimal boundary enclosure $\mathcal H$ whose exterior toward infinity is one ended and has no larger enclosing compact minimal boundary. Moreover,
\[
 \Area_{\widetilde g}(\mathcal H)\ge A_*,
\]
and restricting the comparison fields to that exterior preserves the two total charges.
\end{lemma}

\begin{proof}
The inner boundary $B$ has strictly negative mean curvature for the normal into $X$. Sufficiently distant coordinate spheres have strictly positive outward mean curvature. Apply the trapped-region theorem of Andersson--Metzger with second tensor zero, required inner boundary $B$, and a large outer sphere \cite[Theorem~7.3]{AnderssonMetzger2009}. A strict inner collar makes the weakly trapped union nonempty. Its full marginal frontier is smooth, compact, and minimal. Take the component of its exterior incident on the distant sphere, together with the whole adjacent frontier, and fill bounded pockets on the inner side.

This construction is independent of sufficiently large truncations. Indeed, a weakly trapped bounding surface cannot have a largest-radius point in the region foliated by strictly mean-convex coordinate spheres: the outward mean-curvature comparison at that point is positive. Thus the full trapped union is confined to a fixed compact region. The frontier separates $B$ from infinity and is nonempty.

If a larger compact minimal enclosure existed, its inner side, including the required collar, would be a permissible weakly trapped set, contrary to maximality. This also covers partial coincidence: components already present may be retained while other components are replaced. The comparison is of entire bounding regions. At a touching point of two minimal frontiers ordered from outside, the maximum principle identifies the coincident component.

The end-connected exterior with this boundary is complete up to the boundary. Restoring the inner regions discarded during preparation makes $\mathcal H$ a smooth cut enclosing $S_0$. Therefore
\[
 \Area_{\widetilde g}(\mathcal H)
 \ge\Area_g(\mathcal H)\ge a_g(S_0)=A_*.
\]
The first inequality uses $\widetilde g\ge g$. The closed comparison flux forms have the same integrals over $\mathcal H$ and a distant sphere, by Stokes' Theorem. Hence both total charges are preserved.
\end{proof}

We use the following Riemannian consequence of the charged spacetime
inequality \cite[Theorem~2.3]{ChargedOriginalData2026}.  Let a smooth
connected oriented three-dimensional exterior be complete with its
nonempty compact outermost full minimal boundary included, have one
strongly asymptotically flat end, and satisfy
\[
 R\ge2(|\cE|^2+|\cB|^2),\qquad
 \operatorname{div}\cE=\operatorname{div}\cB=0.
\]
Assume metric decay $O_2(r^{-1})$, field decay $O_1(r^{-2})$, and
integrable scalar curvature.  Then, for ADM mass $\mathcal M$, total
charge magnitude $\mathcal Q$, and full boundary area $4\pi r_H^2$,
\begin{equation}\label{ch:final:kwy}
 \mathcal M\ge\mathcal Q,\qquad
 r_H\le\mathcal M+\sqrt{\mathcal M^2-\mathcal Q^2}.
\end{equation}
The boundary may be disconnected, and no area--charge restriction is
imposed.

Here are the two reductions needed to apply the cited theorem.
If $\mathcal Q>0$, put
\[
 X=\frac{\mathcal Q_E\cE+\mathcal Q_B\cB}{\mathcal Q},
 \qquad
 \cE'=\frac{\mathcal Q_E}{\mathcal Q}X,
 \qquad
 \cB'=\frac{\mathcal Q_B}{\mathcal Q}X.
\]
These fields are divergence free, have the same two total charges and
falloffs, and obey
\[
 \cE'\times\cB'=0,\qquad
 |\cE'|^2+|\cB'|^2=|X|^2\le|\cE|^2+|\cB|^2.
\]
If $\mathcal Q=0$, take both new fields zero.  With second tensor
zero the resulting spacetime data have $J_m=0$ and
$8\pi\mu_m=R/2-(|\cE'|^2+|\cB'|^2)\ge0$, so they satisfy the
matter DEC of the cited theorem.  Their boundary has zero future
expansion.

Secondly, the full enclosing infimum is the full area of an outermost
minimal boundary.  To see this, if a full cut had smaller area, minimize
perimeter among its full enclosing regions containing the original
inner obstacle.  Large mean-convex end spheres confine this problem;
BV compactness and lower semicontinuity give a minimizer.  Its free
frontier is smooth and minimal in dimension three.  At contact with the
smooth minimal obstacle, the minimizing-hull regularity and the strong
comparison principle give smooth coincidence; whole coincident
components are retained.  Filling bounded complementary pockets leaves
the full frontier incident on the end.  Since its area is strictly
smaller, it cannot be just the original boundary.  It is therefore a
strictly larger full minimal enclosure, contrary to outermostness.
This proves the area assertion, with every component counted.
Theorem~2.3 of \cite{ChargedOriginalData2026} now gives
Equation~\eqref{ch:final:kwy}, including positivity of the mass and both
area branches.  Only its numerical conclusion is used.

\begin{proof}[Proof of Theorem~\ref{ch:intro:main}]
Lemma~\ref{ch:rest:timelike} gives a future timelike ADM vector, and
Proposition~\ref{ch:rest:reduction} supplies strict charged rest data
with energy $E_*$ and enclosing infimum $A_*$. Fix one such datum with positive energy and prepare its exterior by Proposition~\ref{ch:prep:prepared}. For each sufficiently small fixed $\epsilon>0$ and sufficiently large $n$, Proposition~\ref{ch:ell:neutral-input} gives a smooth limiting solution after exhausting the truncations at that fixed $n$.

The comparison construction gives the scalar inequality required in \eqref{ch:final:kwy}, closed field forms with the original total charges, and a metric no smaller than the prepared metric. Proposition~\ref{ch:mass:comparison-energy} supplies strong asymptotic flatness, integrable scalar curvature, and an energy bound
\[
 \widetilde E_n\le e^{2\epsilon}(E_*+\eta_n),
 \qquad \eta_n\longrightarrow0
\]
at fixed data and $\epsilon$. Lemma~\ref{ch:final:minimal-enclosure} supplies the appropriate outermost minimal boundary. All hypotheses of the Riemannian theorem are therefore satisfied on its exterior. In particular,
\begin{equation}\label{ch:final:comparison-bound}
 Q\le\widetilde E_n,\qquad
 \sqrt{A_*/(4\pi)}
 \le\widetilde E_n+\sqrt{\widetilde E_n^2-Q^2}.
\end{equation}

Let $F_Q(x)=x+\sqrt{x^2-Q^2}$ for $x\ge Q$. This function is continuous and nondecreasing on its entire domain, including $x=Q$. By \eqref{ch:final:comparison-bound}, the upper energy bound
$e^{2\epsilon}(E_*+\eta_n)$ also lies in that domain. Replacing $\widetilde E_n$ by this bound in \eqref{ch:final:comparison-bound}, then sending $n\to\infty$, gives
\[
 e^{2\epsilon}E_*\ge Q,\qquad
 \sqrt{A_*/(4\pi)}\le F_Q(e^{2\epsilon}E_*).
\]
Letting $\epsilon\downarrow0$ yields
\begin{equation}\label{ch:final:strict-data-bound}
 E_*\ge Q,\qquad
 \sqrt{A_*/(4\pi)}\le E_*+\sqrt{E_*^2-Q^2}.
\end{equation}
Only numerical upper bounds have been passed to the limit; no limiting comparison metric as $n\to\infty$ is needed.

Finally remove the strict rest-data approximation. Proposition~\ref{ch:rest:reduction} keeps $Q_E,Q_B$ fixed, makes $E_*\to m=\sqrt{E^2-|P|^2}$, and gives $A_*\to a_g(S)$. Continuity in \eqref{ch:final:strict-data-bound} proves the theorem. This order of limits requires no constants uniform across the strict-data sequence.
\end{proof}

\begin{remark}[The area branches]
For $0<r\le Q$, the conclusion follows from $m\ge Q$. For $r>Q$, it is equivalent to
\[
 m\ge\frac12\left(r+\frac{Q^2}{r}\right).
\]
The proof uses the upper-area formulation throughout, so disconnected boundaries of either type are covered. When $Q=0$, it reduces to the neutral bound for $a_g(S)$;
Corollary~\ref{bridge:bridge:timelike}, applied as in
Lemma~\ref{ch:rest:timelike}, includes the non-timelike endpoint exclusion.
\end{remark}
